% !TeX root = Thesis.tex

% ****************************************************************************************************
% classicthesis-config.tex
% formerly known as loadpackages.sty, classicthesis-ldpkg.sty, and classicthesis-preamble.sty
% Use it at the beginning of your ClassicThesis.tex, or as a LaTeX Preamble
% in your ClassicThesis.{tex,lyx} with % !TeX root = Thesis.tex

% ****************************************************************************************************
% classicthesis-config.tex
% formerly known as loadpackages.sty, classicthesis-ldpkg.sty, and classicthesis-preamble.sty
% Use it at the beginning of your ClassicThesis.tex, or as a LaTeX Preamble
% in your ClassicThesis.{tex,lyx} with % !TeX root = Thesis.tex

% ****************************************************************************************************
% classicthesis-config.tex
% formerly known as loadpackages.sty, classicthesis-ldpkg.sty, and classicthesis-preamble.sty
% Use it at the beginning of your ClassicThesis.tex, or as a LaTeX Preamble
% in your ClassicThesis.{tex,lyx} with % !TeX root = Thesis.tex

% ****************************************************************************************************
% classicthesis-config.tex
% formerly known as loadpackages.sty, classicthesis-ldpkg.sty, and classicthesis-preamble.sty
% Use it at the beginning of your ClassicThesis.tex, or as a LaTeX Preamble
% in your ClassicThesis.{tex,lyx} with \input{classicthesis-config}
% ****************************************************************************************************
% If you like the classicthesis, then I would appreciate a postcard.
% My address can be found in the file ClassicThesis.pdf. A collection
% of the postcards I received so far is available online at
% http://postcards.miede.de
% ****************************************************************************************************

% set important options before any package can load any of the packages below
\PassOptionsToPackage{dvipsnames}{xcolor} % for classicthesis (e.g. pass before loading tikz)
\PassOptionsToPackage{hyperfootnotes=false,pdfa}{hyperref} % for classicthesis and pdfx

% ****************************************************************************************************
% 0. Set the encoding of your files. UTF-8 is the only sensible encoding nowadays. If you can't read
% äöüßáéçèê∂åëæƒÏ€ then change the encoding setting in your editor, not the line below. If your editor
% does not support utf8 use another editor!
% ****************************************************************************************************
\usepackage[utf8]{inputenc}
\usepackage[T1]{fontenc} % T2A for cyrillics


% ****************************************************************************************************
% 1. Configure classicthesis for your needs here, e.g., remove "drafting" below
% in order to deactivate the time-stamp on the pages
% ****************************************************************************************************
\PassOptionsToPackage{
	% -- TemplateKnob
	drafting=true,            % print version information on the bottom of the pages
	tocaligned=false,         % the left column of the toc will be aligned (no indentation)
	dottedtoc=true,           % page numbers in ToC flushed right
	eulerchapternumbers=true, % use AMS Euler for chapter font (otherwise Palatino)
	floatperchapter=false,    % numbering per chapter for all floats (i.e., Figure 1.1)
	eulermath=true,           % use awesome Euler fonts for mathematical formulae (only with pdfLaTeX)
	beramono=true,            % toggle a nice monospaced font (w/ bold)
	palatino=false,            % deactivate standard font for loading another one, see the last section at the end of this file for suggestions
	%linedheaders=true,       % obsolete / available for backwards compatibility
	style=classicthesis       % classicthesis, arsclassica, linedheaders, plain
}{classicthesis}

% classicthesis overrides for special Makefile test targets
\ifdefined\injectCToptions
	\injectCToptions
\fi

% ****************************************************************************************************
% 2. Language, time and dates, personal data and custom classicthesis template modifications
% ****************************************************************************************************
\usepackage[ngerman,american]{babel} % change this to your language(s), main language last
% Spanish languages need extra options in order to work with this template
%\usepackage[spanish,es-lcroman]{babel}
\usepackage[useregional]{datetime2} % for formatting dates (load after babel/polyglossia)
\usepackage{etoolbox}
% load personal information like title, name, dates, ...
\input{00-definitions/PersonalInfo}

\PassOptionsToPackage{
	% -- TemplateKnob (also have a look into classicthesis-seemoo.sty for further options)
	layout=standard,      % pick from standard, balanced, tight and dense (not recommended)
	% layout definitions:
	%   - standard: the default classicthesis layout with big margins (~69 CPL, 46 lines)
	%   - balanced: wider text area, but still useful margins (~75 CPL, 42 lines)
	%   - tight: even wider text area, tight margin, BCOR > 5 mm looks bad (~77 CPL, 41 lines)
	%   - dense: tiny margins, BCOR > 5 mm almost impossible, not recommended (~87 CPL, 46 lines)
	% how to pick:
	%   - I want to use margin notes --> standard or balanced (tight can be ok with BCOR <= 5 mm)
	%   - I value readability (<80 CPL, up to ~500 words per page) --> standard, balanced, tight
	%   - I have lots of plots and tables --> balanced, tight, dense
	%   - I need as much text as possible per page --> dense
	% notes:
	%   - standard, balanced and tight are roughly equal in the amount of information per page
	%   - balanced and tight are optimized for dina4 with 11 pt Palatino font
	osf=false,            % enable old style figure fonts (not recommended)
	cleantoc=false,       % remove less important chapters from the ToC
	%hidefrontbackmatter, % hide the front and back matter PDF bookmarks
	titleiscover=true,    % treat the title page as the cover (i.e. center the title page)
	%parts,               % use parts (as in a book; for long theses only!)
	%phd,                 % use for a full-blown PhD thesis (with parts, author references, etc.)
	%version=0.1.0,       % when drafting: the version number is displayed on the bottom of the pages
}{classicthesis-seemoo}

% classicthesis-seemoo overrides for special Makefile test targets
\ifdefined\injectCTSoptions
	\injectCTSoptions
\fi
% ****************************************************************************************************


% ****************************************************************************************************
% 3. Loading some handy packages
% ****************************************************************************************************

% ********************************************************************
% BibLaTeX for compiling bibliographies
% ********************************************************************
\usepackage[%
	backend=biber,bibencoding=utf8, % instead of bibtex
	sorting=snyt,                   % shorthand, name, year, title
	maxbibnames=10,                 % default: 3, et al.
	defernumbers=true,              % enable so split references (author's publications) have continuous numbers
	natbib=true,                    % natbib compatibility mode (\citep and \citet still work)
	language=auto,
	style=numeric-comp
]{biblatex}
\DeclareSortingTemplate{snyt}{
	\sort{
		\field{shorthand}
	}
	\sort{
		\field{author}
	}
	\sort{
		\field{year}
	}
	\sort{
		\field{title}
	}
}

% ********************************************************************
% General useful packages
% ********************************************************************
\usepackage[fleqn]{amsmath} % math environments and more by the AMS (fleqn=flush left equations)
\usepackage{graphicx}       % include figures (jpg, png, pdf, ...)
\usepackage{scrhack}        % fix warnings when using KOMA with listings package
\usepackage{xspace}         % to get the spacing after macros right

% optional packages for certain document variants or specific problems
\usepackage[autostyle=true]{csquotes} % for citing in the PhD statement
%\usepackage[LabelsAligned,NoDate]{currvita} % nice cv style for PhD theses (required for the CV)
\usepackage{pdfpages}                % include papers, e.g., for cumulative PhD theses
%\usepackage{xurl}                    % break URLs within word boundaries (not recommended)

% useful packages for drafting
\usepackage{lipsum}         % for generating template text
\usepackage{todonotes}      % manage todos with todonotes

% ********************************************************************
% TikZ and PGFPlots
% ********************************************************************
\usepackage{tikz}
\usepackage{pgfplots}
\pgfplotsset{compat=newest}
\usetikzlibrary{
	chains,
	positioning,
	quotes,
}

\def\StripPrefix#1>{}
\def\isOverleaf{\fi
	\def\overleafJobname{output}% overleaf defaults to 'output' as \jobname
	\edef\overleafJobname{\expandafter\StripPrefix\meaning\overleafJobname}%
	\edef\job{\jobname}%
	\ifx\job\overleafJobname
		}

		% To cache tikz pictures you have to run pdflatex with -shell-escape or --enable-write18
		\ifnum\pdfshellescape=1
			\usepgfplotslibrary{external}
			\if\isOverleaf
				\tikzexternaldisable
			\else
				\tikzexternalize
			\fi
		\fi

		% Lengths for matlab2tikz
		\newlength\figureheight
		\newlength\figurewidth
		% ****************************************************************************************************


		% ****************************************************************************************************
		% 4. Setup floats: tables, (sub)figures, and captions
		% ****************************************************************************************************
		\graphicspath{{06-pictures/}} %path for figures
		\usepackage{float}
		\usepackage{tabularx} % better tables
		\usepackage{booktabs}   % professionelle Tabellenlinien
		\usepackage{longtable}  % mehrseitige Tabellen
		\usepackage{pdflscape}  % Querformat für einzelne Seiten
		\setlength{\extrarowheight}{3pt} % increase table row height
		\usepackage{multirow} % for table cells spanning multiple rows
		\usepackage{subcaption} % for subfigures, subtables, ... (subcaption supersedes subfig)
		\usepackage{makecell}
		% ****************************************************************************************************


		% ****************************************************************************************************
		% 5. Setup code listings
		% ****************************************************************************************************
		\usepackage{listings}
		\lstset{
		% color scheme follows template
		commentstyle=\color{CTsemi},
		keywordstyle={\color{CTtitle}},
		stringstyle=\color{CTcitation},
		basicstyle=\ttfamily\lst@ifdisplaystyle\small\fi, % use normal font size in \lstinline
	emphstyle={\color{CTlink}},
	tabsize=2,
	showstringspaces=false,
	captionpos=b, % caption below listing
	breaklines=true,
	frame=tb,
	numberstyle=\scriptsize,
	numbers=left,
	stepnumber=1,
	numbersep=8pt,
}
% ****************************************************************************************************


% ****************************************************************************************************
% 6. Last calls before the bar closes
% ****************************************************************************************************
% ********************************************************************
% Her Majesty herself
% ********************************************************************
\usepackage{classicthesis}
\usepackage[osf]{newpxtext}

% pdfx for PDF/A (must be loaded after classicthesis, otherwise there will be an issue with the PDF
% bookmarks from the second chapter onwards pointing to the previous page. Seems like an interaction
% between koma-script, titlesec (classicthesis) and hyperref (classicthesis and pdfx).
\usepackage{colorprofiles}       % required for pdfx
\usepackage[a-2b, mathxmp]{pdfx} % reads \jobname.xmpdata and expands commands upon being loaded

% load custom modifications to classicthesis and macros from SEEMOO
\usepackage{classicthesis-seemoo}

% ********************************************************************
% Use the glossaries package for acronyms and fix the long style
% Note: load after hyperref, babel, polyglossia, inputenc and fontenc
% ********************************************************************
\usepackage[
	nopostdot,
	acronym,
	shortcuts,
	nonumberlist,
	nolist]{glossaries}
\makeglossaries

% prevent the long glossary style from incrementing the table counter
\newglossarystyle{fixedlong}{
	\setglossarystyle{long}
	\renewenvironment{theglossary}{
		\edef\savedtablecounter{\the\value{table}}
		\begin{longtable}{lp{\glsdescwidth}}
			}{
		\end{longtable}
		\setcounter{table}{\savedtablecounter}
	}
}
\setglossarystyle{fixedlong}

% ********************************************************************
% Fine-tune hyperreferences (hyperref should be called last)
% ********************************************************************
\hypersetup{%
	% -- TemplateKnob
	%draft, % hyperref's draft mode, for printing see below
	colorlinks=true, linktocpage=true,%
	% uncomment the following line if you want to have black links (e.g., for printing)
	%colorlinks=false, linktocpage=false, pdfborder={0 0 0},%
	breaklinks=true, pageanchor=true,%
	% configure bookmark stuff
	plainpages=false, bookmarksnumbered, bookmarksopen=true, bookmarksopenlevel=1,%
	hypertexnames=true,%
	% define link colors with colors from classicthesis
	urlcolor=CTurl, linkcolor=CTlink, citecolor=CTcitation,%
	%urlcolor=Black, linkcolor=Black, citecolor=Black,
	% pdfinfo is already set via pdfx (c.f. the PDF/A workflow)
}


% ********************************************************************
% Setup autoreferences (hyperref and babel)
% ********************************************************************
% There are some issues regarding autorefnames
% http://www.tex.ac.uk/cgi-bin/texfaq2html?label=latexwords
% you have to redefine the macros for the
% language you use, e.g., american, ngerman
% (as chosen when loading babel/AtBeginDocument)
% ********************************************************************
\makeatletter
\@ifpackageloaded{babel}%
{%
	\addto\extrasamerican{%
		\renewcommand*{\figureautorefname}{Figure}%
		\renewcommand*{\tableautorefname}{Table}%
		\renewcommand*{\partautorefname}{Part}%
		\renewcommand*{\chapterautorefname}{Chapter}%
		\renewcommand*{\sectionautorefname}{Section}%
		\renewcommand*{\subsectionautorefname}{Section}%
		\renewcommand*{\subsubsectionautorefname}{Section}%
	}%
	\addto\extrasngerman{%
		\renewcommand*{\paragraphautorefname}{Absatz}%
		\renewcommand*{\subparagraphautorefname}{Unterabsatz}%
		\renewcommand*{\footnoteautorefname}{Fu\"snote}%
		\renewcommand*{\FancyVerbLineautorefname}{Zeile}%
		\renewcommand*{\theoremautorefname}{Theorem}%
		\renewcommand*{\appendixautorefname}{Anhang}%
		\renewcommand*{\equationautorefname}{Gleichung}%
		\renewcommand*{\itemautorefname}{Punkt}%
	}%
	% Fix to getting autorefs for subfigures right (thanks to Belinda Vogt for changing the definition)
	\providecommand{\subfigureautorefname}{\figureautorefname}%
}{\relax}
\makeatother

% ********************************************************************
% (Better) alternative to \autoref is \cref via the cleveref package
% - load at the very last (after hyperref, amsmath, babel, ...)
% ********************************************************************
\usepackage[capitalize, noabbrev, ngerman, english]{cleveref} % need to pass languages explicitly
%\crefformat{part}{Part #2\MakeUppercase{#1}#3}


% ********************************************************************
% Development Stuff
% ********************************************************************
\listfiles

% ****************************************************************************************************
% If you like the classicthesis, then I would appreciate a postcard.
% My address can be found in the file ClassicThesis.pdf. A collection
% of the postcards I received so far is available online at
% http://postcards.miede.de
% ****************************************************************************************************

% set important options before any package can load any of the packages below
\PassOptionsToPackage{dvipsnames}{xcolor} % for classicthesis (e.g. pass before loading tikz)
\PassOptionsToPackage{hyperfootnotes=false,pdfa}{hyperref} % for classicthesis and pdfx

% ****************************************************************************************************
% 0. Set the encoding of your files. UTF-8 is the only sensible encoding nowadays. If you can't read
% äöüßáéçèê∂åëæƒÏ€ then change the encoding setting in your editor, not the line below. If your editor
% does not support utf8 use another editor!
% ****************************************************************************************************
\usepackage[utf8]{inputenc}
\usepackage[T1]{fontenc} % T2A for cyrillics


% ****************************************************************************************************
% 1. Configure classicthesis for your needs here, e.g., remove "drafting" below
% in order to deactivate the time-stamp on the pages
% ****************************************************************************************************
\PassOptionsToPackage{
	% -- TemplateKnob
	drafting=true,            % print version information on the bottom of the pages
	tocaligned=false,         % the left column of the toc will be aligned (no indentation)
	dottedtoc=true,           % page numbers in ToC flushed right
	eulerchapternumbers=true, % use AMS Euler for chapter font (otherwise Palatino)
	floatperchapter=false,    % numbering per chapter for all floats (i.e., Figure 1.1)
	eulermath=true,           % use awesome Euler fonts for mathematical formulae (only with pdfLaTeX)
	beramono=true,            % toggle a nice monospaced font (w/ bold)
	palatino=false,            % deactivate standard font for loading another one, see the last section at the end of this file for suggestions
	%linedheaders=true,       % obsolete / available for backwards compatibility
	style=classicthesis       % classicthesis, arsclassica, linedheaders, plain
}{classicthesis}

% classicthesis overrides for special Makefile test targets
\ifdefined\injectCToptions
	\injectCToptions
\fi

% ****************************************************************************************************
% 2. Language, time and dates, personal data and custom classicthesis template modifications
% ****************************************************************************************************
\usepackage[ngerman,american]{babel} % change this to your language(s), main language last
% Spanish languages need extra options in order to work with this template
%\usepackage[spanish,es-lcroman]{babel}
\usepackage[useregional]{datetime2} % for formatting dates (load after babel/polyglossia)
\usepackage{etoolbox}
% load personal information like title, name, dates, ...
% !TeX root = ../Thesis.tex

% ##################################################################################################
% -- TemplateKnob
%   >> Personal information to be filled in by thesis author !!!   <<
% NOTES:
% - use \xspace to ensure space after a newcommands (but not for \metaMy... macros and URLs)
% - babel must be loaded for bilingual name definitions
% - datetime2 must be loaded before this file is included to store dates via \DTMsavedate
% ##################################################################################################

% --------------------------------------------------------------------------------------------------
% Main information (required for any type of thesis)
% --------------------------------------------------------------------------------------------------
% IMPORTANT: \metaMy... commands must not have \xspace to avoid errors with pdfx
\newcommand{\metaMyTitle}{YOLO-basierte Detektion und Vergleich visueller Merkmale auf antiken Münzbildern}                  % your title
\newcommand{\metaMyName}{Lisa Haußmann und Dominik Schmitt-Klink}                 % your name
\newcommand{\metaMySubject}{Short summary of the contents}    % avoid &, %, #, \xspace etc.

\DTMsavedate{mySubmissionDate}{2026-10-28}                    % yyyy-mm-dd: thesis submission date
\newcommand{\myProf}{Dr. Karsten Tolle \xspace}  % advisor (bsc/msc) or first examiner

% DO NOT CHANGE: meta data for TUbama PDF/A compliance (pdfx includes the .xmpdata file when loaded)
\begin{filecontents*}{\jobname.xmpdata}
	\Title{\metaMyTitle}
	\Author{\metaMyName}
	\Copyright{Copyright \copyright\ \DTMfetchyear{mySubmissionDate}\ "\metaMyName"}
	\Subject{\metaMySubject}
\end{filecontents*}

% DO NOT CHANGE: derived macros
\newcommand{\myTitle}{\metaMyTitle\xspace}
\newcommand{\myName}{\metaMyName\xspace}


% --------------------------------------------------------------------------------------------------
% Further information (Bachelor's or Master's theses only)
% --------------------------------------------------------------------------------------------------
%\newcommand{\myDegree}{\iflanguage{ngerman}{Bachelorarbeit}{Bachelor's Thesis}\xspace}
\newcommand{\myDegree}{\iflanguage{ngerman}{Bachelorarbeit}{Master's Thesis}\xspace}
\newcommand{\mySupervisor}{Sebastian Gampe\xspace}         % usually the name of a WiMi
\newcommand{\myThesiscode}{SEEMOO-MSC-$0000$\xspace}          % you will get this from our secretary


% --------------------------------------------------------------------------------------------------
% General information of SEEMOO and the TU Darmstadt.
% --------------------------------------------------------------------------------------------------
\newcommand{\myFaculty}{Department of Computer Science\xspace}
\newcommand{\myFacultyDE}{Fachbereich Informatik\xspace}
\newcommand{\myDepartment}{Big Data Lab\xspace}
\newcommand{\myDepartmentDE}{Big Data Lab\xspace}
\newcommand{\myUni}{Goethe Universität Frankfurt am Main\xspace}
\newcommand{\myUniKennziffer}{D17\xspace}
\newcommand{\myLocation}{Frankfurt\xspace}


% --------------------------------------------------------------------------------------------------
% Further information (PhD theses only)
% --------------------------------------------------------------------------------------------------
\newcommand{\myURN}{urn:nbn:de:tuda-tuprints-83253}               % DO NOT PUT \xspace here
\newcommand{\myDegreePhD}{Doktor-Ingenieur (Dr.-Ing.)\xspace}
\newcommand{\myOtherProf}{Dr. The Anh Vuong\xspace}                   % second examiner
\DTMsavedate{myDisputationDate}{1984-12-31}                       % yyyy-mm-dd: disputation date

% the following are only relevant if you include the CV
\DTMsavedate{myBirthDate}{1970-01-01}                             % yyyy-mm-dd: birth date
\newcommand{\myBirthPlace}{Darmstadt, Deutschland\xspace}
\newcommand{\myNationality}{German\xspace}


\PassOptionsToPackage{
	% -- TemplateKnob (also have a look into classicthesis-seemoo.sty for further options)
	layout=standard,      % pick from standard, balanced, tight and dense (not recommended)
	% layout definitions:
	%   - standard: the default classicthesis layout with big margins (~69 CPL, 46 lines)
	%   - balanced: wider text area, but still useful margins (~75 CPL, 42 lines)
	%   - tight: even wider text area, tight margin, BCOR > 5 mm looks bad (~77 CPL, 41 lines)
	%   - dense: tiny margins, BCOR > 5 mm almost impossible, not recommended (~87 CPL, 46 lines)
	% how to pick:
	%   - I want to use margin notes --> standard or balanced (tight can be ok with BCOR <= 5 mm)
	%   - I value readability (<80 CPL, up to ~500 words per page) --> standard, balanced, tight
	%   - I have lots of plots and tables --> balanced, tight, dense
	%   - I need as much text as possible per page --> dense
	% notes:
	%   - standard, balanced and tight are roughly equal in the amount of information per page
	%   - balanced and tight are optimized for dina4 with 11 pt Palatino font
	osf=false,            % enable old style figure fonts (not recommended)
	cleantoc=false,       % remove less important chapters from the ToC
	%hidefrontbackmatter, % hide the front and back matter PDF bookmarks
	titleiscover=true,    % treat the title page as the cover (i.e. center the title page)
	%parts,               % use parts (as in a book; for long theses only!)
	%phd,                 % use for a full-blown PhD thesis (with parts, author references, etc.)
	%version=0.1.0,       % when drafting: the version number is displayed on the bottom of the pages
}{classicthesis-seemoo}

% classicthesis-seemoo overrides for special Makefile test targets
\ifdefined\injectCTSoptions
	\injectCTSoptions
\fi
% ****************************************************************************************************


% ****************************************************************************************************
% 3. Loading some handy packages
% ****************************************************************************************************

% ********************************************************************
% BibLaTeX for compiling bibliographies
% ********************************************************************
\usepackage[%
	backend=biber,bibencoding=utf8, % instead of bibtex
	sorting=snyt,                   % shorthand, name, year, title
	maxbibnames=10,                 % default: 3, et al.
	defernumbers=true,              % enable so split references (author's publications) have continuous numbers
	natbib=true,                    % natbib compatibility mode (\citep and \citet still work)
	language=auto,
	style=numeric-comp
]{biblatex}
\DeclareSortingTemplate{snyt}{
	\sort{
		\field{shorthand}
	}
	\sort{
		\field{author}
	}
	\sort{
		\field{year}
	}
	\sort{
		\field{title}
	}
}

% ********************************************************************
% General useful packages
% ********************************************************************
\usepackage[fleqn]{amsmath} % math environments and more by the AMS (fleqn=flush left equations)
\usepackage{graphicx}       % include figures (jpg, png, pdf, ...)
\usepackage{scrhack}        % fix warnings when using KOMA with listings package
\usepackage{xspace}         % to get the spacing after macros right

% optional packages for certain document variants or specific problems
\usepackage[autostyle=true]{csquotes} % for citing in the PhD statement
%\usepackage[LabelsAligned,NoDate]{currvita} % nice cv style for PhD theses (required for the CV)
\usepackage{pdfpages}                % include papers, e.g., for cumulative PhD theses
%\usepackage{xurl}                    % break URLs within word boundaries (not recommended)

% useful packages for drafting
\usepackage{lipsum}         % for generating template text
\usepackage{todonotes}      % manage todos with todonotes

% ********************************************************************
% TikZ and PGFPlots
% ********************************************************************
\usepackage{tikz}
\usepackage{pgfplots}
\pgfplotsset{compat=newest}
\usetikzlibrary{
	chains,
	positioning,
	quotes,
}

\def\StripPrefix#1>{}
\def\isOverleaf{\fi
	\def\overleafJobname{output}% overleaf defaults to 'output' as \jobname
	\edef\overleafJobname{\expandafter\StripPrefix\meaning\overleafJobname}%
	\edef\job{\jobname}%
	\ifx\job\overleafJobname
		}

		% To cache tikz pictures you have to run pdflatex with -shell-escape or --enable-write18
		\ifnum\pdfshellescape=1
			\usepgfplotslibrary{external}
			\if\isOverleaf
				\tikzexternaldisable
			\else
				\tikzexternalize
			\fi
		\fi

		% Lengths for matlab2tikz
		\newlength\figureheight
		\newlength\figurewidth
		% ****************************************************************************************************


		% ****************************************************************************************************
		% 4. Setup floats: tables, (sub)figures, and captions
		% ****************************************************************************************************
		\graphicspath{{06-pictures/}} %path for figures
		\usepackage{float}
		\usepackage{tabularx} % better tables
		\usepackage{booktabs}   % professionelle Tabellenlinien
		\usepackage{longtable}  % mehrseitige Tabellen
		\usepackage{pdflscape}  % Querformat für einzelne Seiten
		\setlength{\extrarowheight}{3pt} % increase table row height
		\usepackage{multirow} % for table cells spanning multiple rows
		\usepackage{subcaption} % for subfigures, subtables, ... (subcaption supersedes subfig)
		\usepackage{makecell}
		% ****************************************************************************************************


		% ****************************************************************************************************
		% 5. Setup code listings
		% ****************************************************************************************************
		\usepackage{listings}
		\lstset{
		% color scheme follows template
		commentstyle=\color{CTsemi},
		keywordstyle={\color{CTtitle}},
		stringstyle=\color{CTcitation},
		basicstyle=\ttfamily\lst@ifdisplaystyle\small\fi, % use normal font size in \lstinline
	emphstyle={\color{CTlink}},
	tabsize=2,
	showstringspaces=false,
	captionpos=b, % caption below listing
	breaklines=true,
	frame=tb,
	numberstyle=\scriptsize,
	numbers=left,
	stepnumber=1,
	numbersep=8pt,
}
% ****************************************************************************************************


% ****************************************************************************************************
% 6. Last calls before the bar closes
% ****************************************************************************************************
% ********************************************************************
% Her Majesty herself
% ********************************************************************
\usepackage{classicthesis}
\usepackage[osf]{newpxtext}

% pdfx for PDF/A (must be loaded after classicthesis, otherwise there will be an issue with the PDF
% bookmarks from the second chapter onwards pointing to the previous page. Seems like an interaction
% between koma-script, titlesec (classicthesis) and hyperref (classicthesis and pdfx).
\usepackage{colorprofiles}       % required for pdfx
\usepackage[a-2b, mathxmp]{pdfx} % reads \jobname.xmpdata and expands commands upon being loaded

% load custom modifications to classicthesis and macros from SEEMOO
\usepackage{classicthesis-seemoo}

% ********************************************************************
% Use the glossaries package for acronyms and fix the long style
% Note: load after hyperref, babel, polyglossia, inputenc and fontenc
% ********************************************************************
\usepackage[
	nopostdot,
	acronym,
	shortcuts,
	nonumberlist,
	nolist]{glossaries}
\makeglossaries

% prevent the long glossary style from incrementing the table counter
\newglossarystyle{fixedlong}{
	\setglossarystyle{long}
	\renewenvironment{theglossary}{
		\edef\savedtablecounter{\the\value{table}}
		\begin{longtable}{lp{\glsdescwidth}}
			}{
		\end{longtable}
		\setcounter{table}{\savedtablecounter}
	}
}
\setglossarystyle{fixedlong}

% ********************************************************************
% Fine-tune hyperreferences (hyperref should be called last)
% ********************************************************************
\hypersetup{%
	% -- TemplateKnob
	%draft, % hyperref's draft mode, for printing see below
	colorlinks=true, linktocpage=true,%
	% uncomment the following line if you want to have black links (e.g., for printing)
	%colorlinks=false, linktocpage=false, pdfborder={0 0 0},%
	breaklinks=true, pageanchor=true,%
	% configure bookmark stuff
	plainpages=false, bookmarksnumbered, bookmarksopen=true, bookmarksopenlevel=1,%
	hypertexnames=true,%
	% define link colors with colors from classicthesis
	urlcolor=CTurl, linkcolor=CTlink, citecolor=CTcitation,%
	%urlcolor=Black, linkcolor=Black, citecolor=Black,
	% pdfinfo is already set via pdfx (c.f. the PDF/A workflow)
}


% ********************************************************************
% Setup autoreferences (hyperref and babel)
% ********************************************************************
% There are some issues regarding autorefnames
% http://www.tex.ac.uk/cgi-bin/texfaq2html?label=latexwords
% you have to redefine the macros for the
% language you use, e.g., american, ngerman
% (as chosen when loading babel/AtBeginDocument)
% ********************************************************************
\makeatletter
\@ifpackageloaded{babel}%
{%
	\addto\extrasamerican{%
		\renewcommand*{\figureautorefname}{Figure}%
		\renewcommand*{\tableautorefname}{Table}%
		\renewcommand*{\partautorefname}{Part}%
		\renewcommand*{\chapterautorefname}{Chapter}%
		\renewcommand*{\sectionautorefname}{Section}%
		\renewcommand*{\subsectionautorefname}{Section}%
		\renewcommand*{\subsubsectionautorefname}{Section}%
	}%
	\addto\extrasngerman{%
		\renewcommand*{\paragraphautorefname}{Absatz}%
		\renewcommand*{\subparagraphautorefname}{Unterabsatz}%
		\renewcommand*{\footnoteautorefname}{Fu\"snote}%
		\renewcommand*{\FancyVerbLineautorefname}{Zeile}%
		\renewcommand*{\theoremautorefname}{Theorem}%
		\renewcommand*{\appendixautorefname}{Anhang}%
		\renewcommand*{\equationautorefname}{Gleichung}%
		\renewcommand*{\itemautorefname}{Punkt}%
	}%
	% Fix to getting autorefs for subfigures right (thanks to Belinda Vogt for changing the definition)
	\providecommand{\subfigureautorefname}{\figureautorefname}%
}{\relax}
\makeatother

% ********************************************************************
% (Better) alternative to \autoref is \cref via the cleveref package
% - load at the very last (after hyperref, amsmath, babel, ...)
% ********************************************************************
\usepackage[capitalize, noabbrev, ngerman, english]{cleveref} % need to pass languages explicitly
%\crefformat{part}{Part #2\MakeUppercase{#1}#3}


% ********************************************************************
% Development Stuff
% ********************************************************************
\listfiles

% ****************************************************************************************************
% If you like the classicthesis, then I would appreciate a postcard.
% My address can be found in the file ClassicThesis.pdf. A collection
% of the postcards I received so far is available online at
% http://postcards.miede.de
% ****************************************************************************************************

% set important options before any package can load any of the packages below
\PassOptionsToPackage{dvipsnames}{xcolor} % for classicthesis (e.g. pass before loading tikz)
\PassOptionsToPackage{hyperfootnotes=false,pdfa}{hyperref} % for classicthesis and pdfx

% ****************************************************************************************************
% 0. Set the encoding of your files. UTF-8 is the only sensible encoding nowadays. If you can't read
% äöüßáéçèê∂åëæƒÏ€ then change the encoding setting in your editor, not the line below. If your editor
% does not support utf8 use another editor!
% ****************************************************************************************************
\usepackage[utf8]{inputenc}
\usepackage[T1]{fontenc} % T2A for cyrillics


% ****************************************************************************************************
% 1. Configure classicthesis for your needs here, e.g., remove "drafting" below
% in order to deactivate the time-stamp on the pages
% ****************************************************************************************************
\PassOptionsToPackage{
	% -- TemplateKnob
	drafting=true,            % print version information on the bottom of the pages
	tocaligned=false,         % the left column of the toc will be aligned (no indentation)
	dottedtoc=true,           % page numbers in ToC flushed right
	eulerchapternumbers=true, % use AMS Euler for chapter font (otherwise Palatino)
	floatperchapter=false,    % numbering per chapter for all floats (i.e., Figure 1.1)
	eulermath=true,           % use awesome Euler fonts for mathematical formulae (only with pdfLaTeX)
	beramono=true,            % toggle a nice monospaced font (w/ bold)
	palatino=false,            % deactivate standard font for loading another one, see the last section at the end of this file for suggestions
	%linedheaders=true,       % obsolete / available for backwards compatibility
	style=classicthesis       % classicthesis, arsclassica, linedheaders, plain
}{classicthesis}

% classicthesis overrides for special Makefile test targets
\ifdefined\injectCToptions
	\injectCToptions
\fi

% ****************************************************************************************************
% 2. Language, time and dates, personal data and custom classicthesis template modifications
% ****************************************************************************************************
\usepackage[ngerman,american]{babel} % change this to your language(s), main language last
% Spanish languages need extra options in order to work with this template
%\usepackage[spanish,es-lcroman]{babel}
\usepackage[useregional]{datetime2} % for formatting dates (load after babel/polyglossia)
\usepackage{etoolbox}
% load personal information like title, name, dates, ...
% !TeX root = ../Thesis.tex

% ##################################################################################################
% -- TemplateKnob
%   >> Personal information to be filled in by thesis author !!!   <<
% NOTES:
% - use \xspace to ensure space after a newcommands (but not for \metaMy... macros and URLs)
% - babel must be loaded for bilingual name definitions
% - datetime2 must be loaded before this file is included to store dates via \DTMsavedate
% ##################################################################################################

% --------------------------------------------------------------------------------------------------
% Main information (required for any type of thesis)
% --------------------------------------------------------------------------------------------------
% IMPORTANT: \metaMy... commands must not have \xspace to avoid errors with pdfx
\newcommand{\metaMyTitle}{YOLO-basierte Detektion und Vergleich visueller Merkmale auf antiken Münzbildern}                  % your title
\newcommand{\metaMyName}{Lisa Haußmann und Dominik Schmitt-Klink}                 % your name
\newcommand{\metaMySubject}{Short summary of the contents}    % avoid &, %, #, \xspace etc.

\DTMsavedate{mySubmissionDate}{2026-10-28}                    % yyyy-mm-dd: thesis submission date
\newcommand{\myProf}{Dr. Karsten Tolle \xspace}  % advisor (bsc/msc) or first examiner

% DO NOT CHANGE: meta data for TUbama PDF/A compliance (pdfx includes the .xmpdata file when loaded)
\begin{filecontents*}{\jobname.xmpdata}
	\Title{\metaMyTitle}
	\Author{\metaMyName}
	\Copyright{Copyright \copyright\ \DTMfetchyear{mySubmissionDate}\ "\metaMyName"}
	\Subject{\metaMySubject}
\end{filecontents*}

% DO NOT CHANGE: derived macros
\newcommand{\myTitle}{\metaMyTitle\xspace}
\newcommand{\myName}{\metaMyName\xspace}


% --------------------------------------------------------------------------------------------------
% Further information (Bachelor's or Master's theses only)
% --------------------------------------------------------------------------------------------------
%\newcommand{\myDegree}{\iflanguage{ngerman}{Bachelorarbeit}{Bachelor's Thesis}\xspace}
\newcommand{\myDegree}{\iflanguage{ngerman}{Bachelorarbeit}{Master's Thesis}\xspace}
\newcommand{\mySupervisor}{Sebastian Gampe\xspace}         % usually the name of a WiMi
\newcommand{\myThesiscode}{SEEMOO-MSC-$0000$\xspace}          % you will get this from our secretary


% --------------------------------------------------------------------------------------------------
% General information of SEEMOO and the TU Darmstadt.
% --------------------------------------------------------------------------------------------------
\newcommand{\myFaculty}{Department of Computer Science\xspace}
\newcommand{\myFacultyDE}{Fachbereich Informatik\xspace}
\newcommand{\myDepartment}{Big Data Lab\xspace}
\newcommand{\myDepartmentDE}{Big Data Lab\xspace}
\newcommand{\myUni}{Goethe Universität Frankfurt am Main\xspace}
\newcommand{\myUniKennziffer}{D17\xspace}
\newcommand{\myLocation}{Frankfurt\xspace}


% --------------------------------------------------------------------------------------------------
% Further information (PhD theses only)
% --------------------------------------------------------------------------------------------------
\newcommand{\myURN}{urn:nbn:de:tuda-tuprints-83253}               % DO NOT PUT \xspace here
\newcommand{\myDegreePhD}{Doktor-Ingenieur (Dr.-Ing.)\xspace}
\newcommand{\myOtherProf}{Dr. The Anh Vuong\xspace}                   % second examiner
\DTMsavedate{myDisputationDate}{1984-12-31}                       % yyyy-mm-dd: disputation date

% the following are only relevant if you include the CV
\DTMsavedate{myBirthDate}{1970-01-01}                             % yyyy-mm-dd: birth date
\newcommand{\myBirthPlace}{Darmstadt, Deutschland\xspace}
\newcommand{\myNationality}{German\xspace}


\PassOptionsToPackage{
	% -- TemplateKnob (also have a look into classicthesis-seemoo.sty for further options)
	layout=standard,      % pick from standard, balanced, tight and dense (not recommended)
	% layout definitions:
	%   - standard: the default classicthesis layout with big margins (~69 CPL, 46 lines)
	%   - balanced: wider text area, but still useful margins (~75 CPL, 42 lines)
	%   - tight: even wider text area, tight margin, BCOR > 5 mm looks bad (~77 CPL, 41 lines)
	%   - dense: tiny margins, BCOR > 5 mm almost impossible, not recommended (~87 CPL, 46 lines)
	% how to pick:
	%   - I want to use margin notes --> standard or balanced (tight can be ok with BCOR <= 5 mm)
	%   - I value readability (<80 CPL, up to ~500 words per page) --> standard, balanced, tight
	%   - I have lots of plots and tables --> balanced, tight, dense
	%   - I need as much text as possible per page --> dense
	% notes:
	%   - standard, balanced and tight are roughly equal in the amount of information per page
	%   - balanced and tight are optimized for dina4 with 11 pt Palatino font
	osf=false,            % enable old style figure fonts (not recommended)
	cleantoc=false,       % remove less important chapters from the ToC
	%hidefrontbackmatter, % hide the front and back matter PDF bookmarks
	titleiscover=true,    % treat the title page as the cover (i.e. center the title page)
	%parts,               % use parts (as in a book; for long theses only!)
	%phd,                 % use for a full-blown PhD thesis (with parts, author references, etc.)
	%version=0.1.0,       % when drafting: the version number is displayed on the bottom of the pages
}{classicthesis-seemoo}

% classicthesis-seemoo overrides for special Makefile test targets
\ifdefined\injectCTSoptions
	\injectCTSoptions
\fi
% ****************************************************************************************************


% ****************************************************************************************************
% 3. Loading some handy packages
% ****************************************************************************************************

% ********************************************************************
% BibLaTeX for compiling bibliographies
% ********************************************************************
\usepackage[%
	backend=biber,bibencoding=utf8, % instead of bibtex
	sorting=snyt,                   % shorthand, name, year, title
	maxbibnames=10,                 % default: 3, et al.
	defernumbers=true,              % enable so split references (author's publications) have continuous numbers
	natbib=true,                    % natbib compatibility mode (\citep and \citet still work)
	language=auto,
	style=numeric-comp
]{biblatex}
\DeclareSortingTemplate{snyt}{
	\sort{
		\field{shorthand}
	}
	\sort{
		\field{author}
	}
	\sort{
		\field{year}
	}
	\sort{
		\field{title}
	}
}

% ********************************************************************
% General useful packages
% ********************************************************************
\usepackage[fleqn]{amsmath} % math environments and more by the AMS (fleqn=flush left equations)
\usepackage{graphicx}       % include figures (jpg, png, pdf, ...)
\usepackage{scrhack}        % fix warnings when using KOMA with listings package
\usepackage{xspace}         % to get the spacing after macros right

% optional packages for certain document variants or specific problems
\usepackage[autostyle=true]{csquotes} % for citing in the PhD statement
%\usepackage[LabelsAligned,NoDate]{currvita} % nice cv style for PhD theses (required for the CV)
\usepackage{pdfpages}                % include papers, e.g., for cumulative PhD theses
%\usepackage{xurl}                    % break URLs within word boundaries (not recommended)

% useful packages for drafting
\usepackage{lipsum}         % for generating template text
\usepackage{todonotes}      % manage todos with todonotes

% ********************************************************************
% TikZ and PGFPlots
% ********************************************************************
\usepackage{tikz}
\usepackage{pgfplots}
\pgfplotsset{compat=newest}
\usetikzlibrary{
	chains,
	positioning,
	quotes,
}

\def\StripPrefix#1>{}
\def\isOverleaf{\fi
	\def\overleafJobname{output}% overleaf defaults to 'output' as \jobname
	\edef\overleafJobname{\expandafter\StripPrefix\meaning\overleafJobname}%
	\edef\job{\jobname}%
	\ifx\job\overleafJobname
		}

		% To cache tikz pictures you have to run pdflatex with -shell-escape or --enable-write18
		\ifnum\pdfshellescape=1
			\usepgfplotslibrary{external}
			\if\isOverleaf
				\tikzexternaldisable
			\else
				\tikzexternalize
			\fi
		\fi

		% Lengths for matlab2tikz
		\newlength\figureheight
		\newlength\figurewidth
		% ****************************************************************************************************


		% ****************************************************************************************************
		% 4. Setup floats: tables, (sub)figures, and captions
		% ****************************************************************************************************
		\graphicspath{{06-pictures/}} %path for figures
		\usepackage{float}
		\usepackage{tabularx} % better tables
		\usepackage{booktabs}   % professionelle Tabellenlinien
		\usepackage{longtable}  % mehrseitige Tabellen
		\usepackage{pdflscape}  % Querformat für einzelne Seiten
		\setlength{\extrarowheight}{3pt} % increase table row height
		\usepackage{multirow} % for table cells spanning multiple rows
		\usepackage{subcaption} % for subfigures, subtables, ... (subcaption supersedes subfig)
		\usepackage{makecell}
		% ****************************************************************************************************


		% ****************************************************************************************************
		% 5. Setup code listings
		% ****************************************************************************************************
		\usepackage{listings}
		\lstset{
		% color scheme follows template
		commentstyle=\color{CTsemi},
		keywordstyle={\color{CTtitle}},
		stringstyle=\color{CTcitation},
		basicstyle=\ttfamily\lst@ifdisplaystyle\small\fi, % use normal font size in \lstinline
	emphstyle={\color{CTlink}},
	tabsize=2,
	showstringspaces=false,
	captionpos=b, % caption below listing
	breaklines=true,
	frame=tb,
	numberstyle=\scriptsize,
	numbers=left,
	stepnumber=1,
	numbersep=8pt,
}
% ****************************************************************************************************


% ****************************************************************************************************
% 6. Last calls before the bar closes
% ****************************************************************************************************
% ********************************************************************
% Her Majesty herself
% ********************************************************************
\usepackage{classicthesis}
\usepackage[osf]{newpxtext}

% pdfx for PDF/A (must be loaded after classicthesis, otherwise there will be an issue with the PDF
% bookmarks from the second chapter onwards pointing to the previous page. Seems like an interaction
% between koma-script, titlesec (classicthesis) and hyperref (classicthesis and pdfx).
\usepackage{colorprofiles}       % required for pdfx
\usepackage[a-2b, mathxmp]{pdfx} % reads \jobname.xmpdata and expands commands upon being loaded

% load custom modifications to classicthesis and macros from SEEMOO
\usepackage{classicthesis-seemoo}

% ********************************************************************
% Use the glossaries package for acronyms and fix the long style
% Note: load after hyperref, babel, polyglossia, inputenc and fontenc
% ********************************************************************
\usepackage[
	nopostdot,
	acronym,
	shortcuts,
	nonumberlist,
	nolist]{glossaries}
\makeglossaries

% prevent the long glossary style from incrementing the table counter
\newglossarystyle{fixedlong}{
	\setglossarystyle{long}
	\renewenvironment{theglossary}{
		\edef\savedtablecounter{\the\value{table}}
		\begin{longtable}{lp{\glsdescwidth}}
			}{
		\end{longtable}
		\setcounter{table}{\savedtablecounter}
	}
}
\setglossarystyle{fixedlong}

% ********************************************************************
% Fine-tune hyperreferences (hyperref should be called last)
% ********************************************************************
\hypersetup{%
	% -- TemplateKnob
	%draft, % hyperref's draft mode, for printing see below
	colorlinks=true, linktocpage=true,%
	% uncomment the following line if you want to have black links (e.g., for printing)
	%colorlinks=false, linktocpage=false, pdfborder={0 0 0},%
	breaklinks=true, pageanchor=true,%
	% configure bookmark stuff
	plainpages=false, bookmarksnumbered, bookmarksopen=true, bookmarksopenlevel=1,%
	hypertexnames=true,%
	% define link colors with colors from classicthesis
	urlcolor=CTurl, linkcolor=CTlink, citecolor=CTcitation,%
	%urlcolor=Black, linkcolor=Black, citecolor=Black,
	% pdfinfo is already set via pdfx (c.f. the PDF/A workflow)
}


% ********************************************************************
% Setup autoreferences (hyperref and babel)
% ********************************************************************
% There are some issues regarding autorefnames
% http://www.tex.ac.uk/cgi-bin/texfaq2html?label=latexwords
% you have to redefine the macros for the
% language you use, e.g., american, ngerman
% (as chosen when loading babel/AtBeginDocument)
% ********************************************************************
\makeatletter
\@ifpackageloaded{babel}%
{%
	\addto\extrasamerican{%
		\renewcommand*{\figureautorefname}{Figure}%
		\renewcommand*{\tableautorefname}{Table}%
		\renewcommand*{\partautorefname}{Part}%
		\renewcommand*{\chapterautorefname}{Chapter}%
		\renewcommand*{\sectionautorefname}{Section}%
		\renewcommand*{\subsectionautorefname}{Section}%
		\renewcommand*{\subsubsectionautorefname}{Section}%
	}%
	\addto\extrasngerman{%
		\renewcommand*{\paragraphautorefname}{Absatz}%
		\renewcommand*{\subparagraphautorefname}{Unterabsatz}%
		\renewcommand*{\footnoteautorefname}{Fu\"snote}%
		\renewcommand*{\FancyVerbLineautorefname}{Zeile}%
		\renewcommand*{\theoremautorefname}{Theorem}%
		\renewcommand*{\appendixautorefname}{Anhang}%
		\renewcommand*{\equationautorefname}{Gleichung}%
		\renewcommand*{\itemautorefname}{Punkt}%
	}%
	% Fix to getting autorefs for subfigures right (thanks to Belinda Vogt for changing the definition)
	\providecommand{\subfigureautorefname}{\figureautorefname}%
}{\relax}
\makeatother

% ********************************************************************
% (Better) alternative to \autoref is \cref via the cleveref package
% - load at the very last (after hyperref, amsmath, babel, ...)
% ********************************************************************
\usepackage[capitalize, noabbrev, ngerman, english]{cleveref} % need to pass languages explicitly
%\crefformat{part}{Part #2\MakeUppercase{#1}#3}


% ********************************************************************
% Development Stuff
% ********************************************************************
\listfiles

% ****************************************************************************************************
% If you like the classicthesis, then I would appreciate a postcard.
% My address can be found in the file ClassicThesis.pdf. A collection
% of the postcards I received so far is available online at
% http://postcards.miede.de
% ****************************************************************************************************

% set important options before any package can load any of the packages below
\PassOptionsToPackage{dvipsnames}{xcolor} % for classicthesis (e.g. pass before loading tikz)
\PassOptionsToPackage{hyperfootnotes=false,pdfa}{hyperref} % for classicthesis and pdfx

% ****************************************************************************************************
% 0. Set the encoding of your files. UTF-8 is the only sensible encoding nowadays. If you can't read
% äöüßáéçèê∂åëæƒÏ€ then change the encoding setting in your editor, not the line below. If your editor
% does not support utf8 use another editor!
% ****************************************************************************************************
\usepackage[utf8]{inputenc}
\usepackage[T1]{fontenc} % T2A for cyrillics


% ****************************************************************************************************
% 1. Configure classicthesis for your needs here, e.g., remove "drafting" below
% in order to deactivate the time-stamp on the pages
% ****************************************************************************************************
\PassOptionsToPackage{
	% -- TemplateKnob
	drafting=true,            % print version information on the bottom of the pages
	tocaligned=false,         % the left column of the toc will be aligned (no indentation)
	dottedtoc=true,           % page numbers in ToC flushed right
	eulerchapternumbers=true, % use AMS Euler for chapter font (otherwise Palatino)
	floatperchapter=false,    % numbering per chapter for all floats (i.e., Figure 1.1)
	eulermath=true,           % use awesome Euler fonts for mathematical formulae (only with pdfLaTeX)
	beramono=true,            % toggle a nice monospaced font (w/ bold)
	palatino=false,            % deactivate standard font for loading another one, see the last section at the end of this file for suggestions
	%linedheaders=true,       % obsolete / available for backwards compatibility
	style=classicthesis       % classicthesis, arsclassica, linedheaders, plain
}{classicthesis}

% classicthesis overrides for special Makefile test targets
\ifdefined\injectCToptions
	\injectCToptions
\fi

% ****************************************************************************************************
% 2. Language, time and dates, personal data and custom classicthesis template modifications
% ****************************************************************************************************
\usepackage[ngerman,american]{babel} % change this to your language(s), main language last
% Spanish languages need extra options in order to work with this template
%\usepackage[spanish,es-lcroman]{babel}
\usepackage[useregional]{datetime2} % for formatting dates (load after babel/polyglossia)
\usepackage{etoolbox}
% load personal information like title, name, dates, ...
% !TeX root = ../Thesis.tex

% ##################################################################################################
% -- TemplateKnob
%   >> Personal information to be filled in by thesis author !!!   <<
% NOTES:
% - use \xspace to ensure space after a newcommands (but not for \metaMy... macros and URLs)
% - babel must be loaded for bilingual name definitions
% - datetime2 must be loaded before this file is included to store dates via \DTMsavedate
% ##################################################################################################

% --------------------------------------------------------------------------------------------------
% Main information (required for any type of thesis)
% --------------------------------------------------------------------------------------------------
% IMPORTANT: \metaMy... commands must not have \xspace to avoid errors with pdfx
\newcommand{\metaMyTitle}{YOLO-basierte Detektion und Vergleich visueller Merkmale auf antiken Münzbildern}                  % your title
\newcommand{\metaMyName}{Lisa Haußmann und Dominik Schmitt-Klink}                 % your name
\newcommand{\metaMySubject}{Short summary of the contents}    % avoid &, %, #, \xspace etc.

\DTMsavedate{mySubmissionDate}{2026-10-28}                    % yyyy-mm-dd: thesis submission date
\newcommand{\myProf}{Dr. Karsten Tolle \xspace}  % advisor (bsc/msc) or first examiner

% DO NOT CHANGE: meta data for TUbama PDF/A compliance (pdfx includes the .xmpdata file when loaded)
\begin{filecontents*}{\jobname.xmpdata}
	\Title{\metaMyTitle}
	\Author{\metaMyName}
	\Copyright{Copyright \copyright\ \DTMfetchyear{mySubmissionDate}\ "\metaMyName"}
	\Subject{\metaMySubject}
\end{filecontents*}

% DO NOT CHANGE: derived macros
\newcommand{\myTitle}{\metaMyTitle\xspace}
\newcommand{\myName}{\metaMyName\xspace}


% --------------------------------------------------------------------------------------------------
% Further information (Bachelor's or Master's theses only)
% --------------------------------------------------------------------------------------------------
%\newcommand{\myDegree}{\iflanguage{ngerman}{Bachelorarbeit}{Bachelor's Thesis}\xspace}
\newcommand{\myDegree}{\iflanguage{ngerman}{Bachelorarbeit}{Master's Thesis}\xspace}
\newcommand{\mySupervisor}{Sebastian Gampe\xspace}         % usually the name of a WiMi
\newcommand{\myThesiscode}{SEEMOO-MSC-$0000$\xspace}          % you will get this from our secretary


% --------------------------------------------------------------------------------------------------
% General information of SEEMOO and the TU Darmstadt.
% --------------------------------------------------------------------------------------------------
\newcommand{\myFaculty}{Department of Computer Science\xspace}
\newcommand{\myFacultyDE}{Fachbereich Informatik\xspace}
\newcommand{\myDepartment}{Big Data Lab\xspace}
\newcommand{\myDepartmentDE}{Big Data Lab\xspace}
\newcommand{\myUni}{Goethe Universität Frankfurt am Main\xspace}
\newcommand{\myUniKennziffer}{D17\xspace}
\newcommand{\myLocation}{Frankfurt\xspace}


% --------------------------------------------------------------------------------------------------
% Further information (PhD theses only)
% --------------------------------------------------------------------------------------------------
\newcommand{\myURN}{urn:nbn:de:tuda-tuprints-83253}               % DO NOT PUT \xspace here
\newcommand{\myDegreePhD}{Doktor-Ingenieur (Dr.-Ing.)\xspace}
\newcommand{\myOtherProf}{Dr. The Anh Vuong\xspace}                   % second examiner
\DTMsavedate{myDisputationDate}{1984-12-31}                       % yyyy-mm-dd: disputation date

% the following are only relevant if you include the CV
\DTMsavedate{myBirthDate}{1970-01-01}                             % yyyy-mm-dd: birth date
\newcommand{\myBirthPlace}{Darmstadt, Deutschland\xspace}
\newcommand{\myNationality}{German\xspace}


\PassOptionsToPackage{
	% -- TemplateKnob (also have a look into classicthesis-seemoo.sty for further options)
	layout=standard,      % pick from standard, balanced, tight and dense (not recommended)
	% layout definitions:
	%   - standard: the default classicthesis layout with big margins (~69 CPL, 46 lines)
	%   - balanced: wider text area, but still useful margins (~75 CPL, 42 lines)
	%   - tight: even wider text area, tight margin, BCOR > 5 mm looks bad (~77 CPL, 41 lines)
	%   - dense: tiny margins, BCOR > 5 mm almost impossible, not recommended (~87 CPL, 46 lines)
	% how to pick:
	%   - I want to use margin notes --> standard or balanced (tight can be ok with BCOR <= 5 mm)
	%   - I value readability (<80 CPL, up to ~500 words per page) --> standard, balanced, tight
	%   - I have lots of plots and tables --> balanced, tight, dense
	%   - I need as much text as possible per page --> dense
	% notes:
	%   - standard, balanced and tight are roughly equal in the amount of information per page
	%   - balanced and tight are optimized for dina4 with 11 pt Palatino font
	osf=false,            % enable old style figure fonts (not recommended)
	cleantoc=false,       % remove less important chapters from the ToC
	%hidefrontbackmatter, % hide the front and back matter PDF bookmarks
	titleiscover=true,    % treat the title page as the cover (i.e. center the title page)
	%parts,               % use parts (as in a book; for long theses only!)
	%phd,                 % use for a full-blown PhD thesis (with parts, author references, etc.)
	%version=0.1.0,       % when drafting: the version number is displayed on the bottom of the pages
}{classicthesis-seemoo}

% classicthesis-seemoo overrides for special Makefile test targets
\ifdefined\injectCTSoptions
	\injectCTSoptions
\fi
% ****************************************************************************************************


% ****************************************************************************************************
% 3. Loading some handy packages
% ****************************************************************************************************

% ********************************************************************
% BibLaTeX for compiling bibliographies
% ********************************************************************
\usepackage[%
	backend=biber,bibencoding=utf8, % instead of bibtex
	sorting=snyt,                   % shorthand, name, year, title
	maxbibnames=10,                 % default: 3, et al.
	defernumbers=true,              % enable so split references (author's publications) have continuous numbers
	natbib=true,                    % natbib compatibility mode (\citep and \citet still work)
	language=auto,
	style=numeric-comp
]{biblatex}
\DeclareSortingTemplate{snyt}{
	\sort{
		\field{shorthand}
	}
	\sort{
		\field{author}
	}
	\sort{
		\field{year}
	}
	\sort{
		\field{title}
	}
}

% ********************************************************************
% General useful packages
% ********************************************************************
\usepackage[fleqn]{amsmath} % math environments and more by the AMS (fleqn=flush left equations)
\usepackage{graphicx}       % include figures (jpg, png, pdf, ...)
\usepackage{scrhack}        % fix warnings when using KOMA with listings package
\usepackage{xspace}         % to get the spacing after macros right

% optional packages for certain document variants or specific problems
\usepackage[autostyle=true]{csquotes} % for citing in the PhD statement
%\usepackage[LabelsAligned,NoDate]{currvita} % nice cv style for PhD theses (required for the CV)
\usepackage{pdfpages}                % include papers, e.g., for cumulative PhD theses
%\usepackage{xurl}                    % break URLs within word boundaries (not recommended)

% useful packages for drafting
\usepackage{lipsum}         % for generating template text
\usepackage{todonotes}      % manage todos with todonotes

% ********************************************************************
% TikZ and PGFPlots
% ********************************************************************
\usepackage{tikz}
\usepackage{pgfplots}
\pgfplotsset{compat=newest}
\usetikzlibrary{
	chains,
	positioning,
	quotes,
}

\def\StripPrefix#1>{}
\def\isOverleaf{\fi
	\def\overleafJobname{output}% overleaf defaults to 'output' as \jobname
	\edef\overleafJobname{\expandafter\StripPrefix\meaning\overleafJobname}%
	\edef\job{\jobname}%
	\ifx\job\overleafJobname
		}

		% To cache tikz pictures you have to run pdflatex with -shell-escape or --enable-write18
		\ifnum\pdfshellescape=1
			\usepgfplotslibrary{external}
			\if\isOverleaf
				\tikzexternaldisable
			\else
				\tikzexternalize
			\fi
		\fi

		% Lengths for matlab2tikz
		\newlength\figureheight
		\newlength\figurewidth
		% ****************************************************************************************************


		% ****************************************************************************************************
		% 4. Setup floats: tables, (sub)figures, and captions
		% ****************************************************************************************************
		\graphicspath{{06-pictures/}} %path for figures
		\usepackage{float}
		\usepackage{tabularx} % better tables
		\usepackage{booktabs}   % professionelle Tabellenlinien
		\usepackage{longtable}  % mehrseitige Tabellen
		\usepackage{pdflscape}  % Querformat für einzelne Seiten
		\setlength{\extrarowheight}{3pt} % increase table row height
		\usepackage{multirow} % for table cells spanning multiple rows
		\usepackage{subcaption} % for subfigures, subtables, ... (subcaption supersedes subfig)
		\usepackage{makecell}
		% ****************************************************************************************************


		% ****************************************************************************************************
		% 5. Setup code listings
		% ****************************************************************************************************
		\usepackage{listings}
		\lstset{
		% color scheme follows template
		commentstyle=\color{CTsemi},
		keywordstyle={\color{CTtitle}},
		stringstyle=\color{CTcitation},
		basicstyle=\ttfamily\lst@ifdisplaystyle\small\fi, % use normal font size in \lstinline
	emphstyle={\color{CTlink}},
	tabsize=2,
	showstringspaces=false,
	captionpos=b, % caption below listing
	breaklines=true,
	frame=tb,
	numberstyle=\scriptsize,
	numbers=left,
	stepnumber=1,
	numbersep=8pt,
}
% ****************************************************************************************************


% ****************************************************************************************************
% 6. Last calls before the bar closes
% ****************************************************************************************************
% ********************************************************************
% Her Majesty herself
% ********************************************************************
\usepackage{classicthesis}
\usepackage[osf]{newpxtext}

% pdfx for PDF/A (must be loaded after classicthesis, otherwise there will be an issue with the PDF
% bookmarks from the second chapter onwards pointing to the previous page. Seems like an interaction
% between koma-script, titlesec (classicthesis) and hyperref (classicthesis and pdfx).
\usepackage{colorprofiles}       % required for pdfx
\usepackage[a-2b, mathxmp]{pdfx} % reads \jobname.xmpdata and expands commands upon being loaded

% load custom modifications to classicthesis and macros from SEEMOO
\usepackage{classicthesis-seemoo}

% ********************************************************************
% Use the glossaries package for acronyms and fix the long style
% Note: load after hyperref, babel, polyglossia, inputenc and fontenc
% ********************************************************************
\usepackage[
	nopostdot,
	acronym,
	shortcuts,
	nonumberlist,
	nolist]{glossaries}
\makeglossaries

% prevent the long glossary style from incrementing the table counter
\newglossarystyle{fixedlong}{
	\setglossarystyle{long}
	\renewenvironment{theglossary}{
		\edef\savedtablecounter{\the\value{table}}
		\begin{longtable}{lp{\glsdescwidth}}
			}{
		\end{longtable}
		\setcounter{table}{\savedtablecounter}
	}
}
\setglossarystyle{fixedlong}

% ********************************************************************
% Fine-tune hyperreferences (hyperref should be called last)
% ********************************************************************
\hypersetup{%
	% -- TemplateKnob
	%draft, % hyperref's draft mode, for printing see below
	colorlinks=true, linktocpage=true,%
	% uncomment the following line if you want to have black links (e.g., for printing)
	%colorlinks=false, linktocpage=false, pdfborder={0 0 0},%
	breaklinks=true, pageanchor=true,%
	% configure bookmark stuff
	plainpages=false, bookmarksnumbered, bookmarksopen=true, bookmarksopenlevel=1,%
	hypertexnames=true,%
	% define link colors with colors from classicthesis
	urlcolor=CTurl, linkcolor=CTlink, citecolor=CTcitation,%
	%urlcolor=Black, linkcolor=Black, citecolor=Black,
	% pdfinfo is already set via pdfx (c.f. the PDF/A workflow)
}


% ********************************************************************
% Setup autoreferences (hyperref and babel)
% ********************************************************************
% There are some issues regarding autorefnames
% http://www.tex.ac.uk/cgi-bin/texfaq2html?label=latexwords
% you have to redefine the macros for the
% language you use, e.g., american, ngerman
% (as chosen when loading babel/AtBeginDocument)
% ********************************************************************
\makeatletter
\@ifpackageloaded{babel}%
{%
	\addto\extrasamerican{%
		\renewcommand*{\figureautorefname}{Figure}%
		\renewcommand*{\tableautorefname}{Table}%
		\renewcommand*{\partautorefname}{Part}%
		\renewcommand*{\chapterautorefname}{Chapter}%
		\renewcommand*{\sectionautorefname}{Section}%
		\renewcommand*{\subsectionautorefname}{Section}%
		\renewcommand*{\subsubsectionautorefname}{Section}%
	}%
	\addto\extrasngerman{%
		\renewcommand*{\paragraphautorefname}{Absatz}%
		\renewcommand*{\subparagraphautorefname}{Unterabsatz}%
		\renewcommand*{\footnoteautorefname}{Fu\"snote}%
		\renewcommand*{\FancyVerbLineautorefname}{Zeile}%
		\renewcommand*{\theoremautorefname}{Theorem}%
		\renewcommand*{\appendixautorefname}{Anhang}%
		\renewcommand*{\equationautorefname}{Gleichung}%
		\renewcommand*{\itemautorefname}{Punkt}%
	}%
	% Fix to getting autorefs for subfigures right (thanks to Belinda Vogt for changing the definition)
	\providecommand{\subfigureautorefname}{\figureautorefname}%
}{\relax}
\makeatother

% ********************************************************************
% (Better) alternative to \autoref is \cref via the cleveref package
% - load at the very last (after hyperref, amsmath, babel, ...)
% ********************************************************************
\usepackage[capitalize, noabbrev, ngerman, english]{cleveref} % need to pass languages explicitly
%\crefformat{part}{Part #2\MakeUppercase{#1}#3}


% ********************************************************************
% Development Stuff
% ********************************************************************
\listfiles
